% =============================================================================
% Statements of the mock midterm (Topics 1, 2A, 2B), written once and \input by
%   mock-midterm-topics-1-2b.tex            (exam paper)
%   mock-midterm-topics-1-2b-solutions.tex  (worked solutions)
% Every statement is a real exercise copied verbatim (rules/exercise-lists.md):
% source line first, then the full text. Allowed edits only: math in LaTeX, tables
% rebuilt, line breaks undone, currency written as \$.
% Needs: \srcline{...} defined by the including document.
% =============================================================================

% ---------------------------------------------------------------- Question 1.1
\newcommand{\QoneA}{%
\srcline{Berk, DeMarzo \& Harford, \emph{Fundamentals of Corporate Finance} 2e, Ch.~2, Exercise 2.24,
p.~55 (PDF 86) $\cdot$ Answer in book: no}

\noindent In January 2009, American Airlines (AMR) had a market capitalization of \$1.7 billion, debt of
\$11.1 billion, and cash of \$4.6 billion. American Airlines had revenues of \$23.8 billion. British
Airways (BABWF) had a market capitalization of \$2.2 billion, debt of \$4.7 billion, cash of \$2.6
billion, and revenues of \$13.1 billion.
\begin{enumerate}[label=\alph*.,leftmargin=2em]
  \item Compare the market capitalization-to-revenue ratio (also called the price-to-sales ratio) for
        American Airlines and British Airways.
  \item Compare the enterprise value-to-revenue ratio for American Airlines and British Airways.
  \item Which of these comparisons is more meaningful? Explain.
\end{enumerate}}

% ---------------------------------------------------------------- Question 1.2
\newcommand{\QoneB}{%
\srcline{Berk, DeMarzo \& Harford, \emph{Fundamentals of Corporate Finance} 2e, Ch.~9, Exercise 9.13,
p.~275 (PDF 306) $\cdot$ Answer in book: no}

\noindent Elmdale Enterprises is deciding whether to expand its production facilities. Although long-term
cash flows are difficult to estimate, management has projected the following cash flows for the first
two years (in millions of dollars):
\begin{center}
\begin{tabular}{lrr}
  \toprule
  Year & 1 & 2 \\ \midrule
  Revenues & 125 & 160 \\
  Operating Expenses (other than depreciation) & 40 & 60 \\
  Depreciation & 25 & 36 \\
  Increase in Net Working Capital & 2 & 8 \\
  Capital Expenditures & 30 & 40 \\
  Marginal Corporate Tax Rate & 35\% & 35\% \\ \bottomrule
\end{tabular}
\end{center}
\begin{enumerate}[label=\alph*.,leftmargin=2em]
  \item What are the incremental earnings for this project for years 1 and 2?
  \item What are the free cash flows for this project for the first two years?
\end{enumerate}}

% ---------------------------------------------------------------- Question 1.3
\newcommand{\QoneC}{%
\srcline{Berk, DeMarzo \& Harford, \emph{Fundamentals of Corporate Finance} 2e, Ch.~16, Exercises
16.19--16.21, pp.~493--494 (PDF 524--525) $\cdot$ Answer in book: no}

\noindent\textbf{19.} Kohwe Corporation plans to issue equity to raise \$50 million to finance a new
investment. After making the investment, Kohwe expects to earn free cash flows of \$10 million each year.
Kohwe currently has 5 million shares outstanding, and it has no other assets or opportunities. Suppose
the appropriate discount rate for Kohwe's future free cash flows is 8\%, and the only capital market
imperfections are corporate taxes and financial distress costs.
\begin{enumerate}[label=\alph*.,leftmargin=2em]
  \item What is the NPV of Kohwe's investment?
  \item What is Kohwe's share price today?
\end{enumerate}
\noindent\textbf{20.} Suppose Kohwe borrows the \$50 million instead. The firm will pay interest only on
this loan each year, and it will maintain an outstanding balance of \$50 million on the loan. Suppose
that Kohwe's corporate tax rate is 40\%, and expected free cash flows are still \$10 million each year.
What is Kohwe's share price today if the investment is financed with debt?

\smallskip\noindent\textbf{21.} Now suppose that with leverage, Kohwe's expected free cash flows will
decline to \$9 million per year due to reduced sales and other financial distress costs. Assume that the
appropriate discount rate for Kohwe's future free cash flows is still 8\%. What is Kohwe's share price
today given the financial distress costs of leverage?}

% ---------------------------------------------------------------- Question 1.4
\newcommand{\QoneD}{%
\srcline{Berk, DeMarzo \& Harford, \emph{Fundamentals of Corporate Finance} 2e, Ch.~17, Exercises
17.11--17.13, p.~528 (PDF 559) $\cdot$ Answer in book: no. Problems 12 and 13 say ``Redo Problem 11'',
so Problem 11 is part of this block.}

\noindent\textbf{11.} Assume perfect capital markets. Kay Industries currently has \$100 million invested
in short-term Treasury securities paying 7\%, and it pays out the interest payments on these securities
as a dividend. The board is considering selling the Treasury securities and paying out the proceeds as a
one-time dividend payment.
\begin{enumerate}[label=\alph*.,leftmargin=2em]
  \item If the board went ahead with this plan, what would happen to the value of Kay stock upon the
        announcement of a change in policy?
  \item What would happen to the value of Kay stock on the ex-dividend date of the one-time dividend?
  \item Given these price reactions, will this decision benefit investors?
\end{enumerate}
\noindent\textbf{12.} Redo Problem 11, but assume that Kay must pay a corporate tax rate of 35\%, and that
investors pay no taxes.

\smallskip\noindent\textbf{13.} Redo Problem 11, but assume that investors pay a 15\% tax on dividends but
no capital gains taxes, and that Kay does not pay corporate taxes.}

% ---------------------------------------------------------------- Question 2 (previous-year midterm)
\newcommand{\QtwoIntro}{%
\srcline{Finance II -- Corporate Finance, \emph{Midterm Exam (Previous Year)}, Question 2, pp.~1--2
(\texttt{Prova/q2.pdf}) $\cdot$ Answer in source: no}

\noindent Company A has no assets and is considering investing in a new project. The expected
\emph{after-tax} unlevered free cash flows are given below:
\begin{center}
\begin{tabular}{lrrrrrr}
  \toprule
  Year & 0 & 1 & 2 & 3 & 4 & 5 \\ \midrule
  FCF & $-750$ & 250 & 250 & 250 & 250 & 250 \\ \bottomrule
\end{tabular}
\end{center}
As the project is new, owners of company A will look at comparable firms to find an appropriate discount
rate. They identify a comparable firm whose assets (\emph{excluding the market value of any tax benefits
from debt}) have the same market or systematic risk as the project that company A is considering
implementing. The equity of the comparable firm trades in financial markets. The market value of the
assets (\emph{excluding the market value of any tax benefits from debt}) of the comparable firm is
denoted by $V_U$ (as if the firm had no debt). If the comparable firm had no leverage, the market value
of its assets would be equal to $V_U$, that is, $A=V_U$. However, as the comparable firm maintains a
constant debt-to-equity ratio of 0.1, the market value of its assets equals $V_U$ plus the present value
of interest tax shields, that is, $A=V_U+PV(ITS)$. The debt of the comparable firm is risk-free and its
equity has a beta of $\beta_E=1.65$. The expected excess market return, $\mathbb{E}[\tilde r_m]-r_f$, is
10\% and the risk free rate is $r_f=5\%$. The assumptions of the CAPM and the Modigliani-Miller model hold
(including competitive financial markets). The corporate tax rate is 30\%.}

\newcommand{\QtwoA}{%
\textbf{(a)} (\emph{10 marks}) What is the beta of the total assets ($\beta_A$), assets excluding tax
benefits of debt ($\beta_U$), debt ($\beta_D$), and tax shields ($\beta_{ITS}$) of the comparable firm?
What is the expected return on the assets ($r_A$), assets excluding tax benefits ($r_U$), debt ($r_D$),
tax shields ($r_{ITS}$), and equity ($r_E$) of the comparable firm?\\
\emph{Hint: we have that}
\[ r_A=\frac DA r_D+\frac EA r_E=\frac{V_U}{A}r_U+\frac{PV(ITS)}{A}r_{ITS} \]
\emph{and}
\[ \beta_A=\frac DA\beta_D+\frac EA\beta_E=\frac{V_U}{A}\beta_U+\frac{PV(ITS)}{A}\beta_{ITS} \]
\emph{but we made an assumption about the capital structure of the comparable firm -- constant
debt-to-equity ratio -- that simplifies this formula.}}

\newcommand{\QtwoB}{%
\textbf{(b)} (\emph{10 marks}) Suppose the owners of company A decide to finance the investment using
equity only. What is the expected return on the assets ($r_A$) and equity ($r_E$) of company A?
Calculate the NPV of the project. Assuming the project is undertaken, what is the market value of the
assets of the firm ($V_U$)? Are outside investors willing to provide financing? If yes, what share of the
firm they buy in exchange for providing 750?}

\newcommand{\QtwoC}{%
\textbf{(c)} (\emph{10 marks}) Suppose the owners of company A decide to maintain a constant
debt-to-equity ratio of 0.5. Debt is still risk-free. What is the required return on equity, $r_E$? Is
it smaller, equal, or larger than the expected return on equity found in (b)? Explain.}

\newcommand{\QtwoD}{%
\textbf{(d)} (\emph{10 marks}) Keeping the assumptions of part (c) -- constant debt-to-equity ratio of 0.5
and risk-free debt --, what would be the NPV of the project? Explain the reason for your findings. Are
investors willing to provide financing?}

\newcommand{\QtwoE}{%
\textbf{(e)} (\emph{10 marks}) Keeping the assumptions of part (c) -- constant debt-to-equity ratio of 0.5
and risk-free debt --, what is the market value of debt (D), equity (E), and the total assets of the firm
at time 0? What is the share of equity that original owners retain and its market value? The value of
the position of original owners is smaller, equal, or larger than the NPV of the project? Explain.}

\newcommand{\QtwoF}{%
\textbf{(f)} (\emph{10 marks}) Keeping the assumptions of part (c) --constant debt-to-equity ratio of 0.5
and risk-free debt --, what is the market value of debt at $t=1$, $t=2$, $t=3$, $t=4$, and $t=5$? Use
these values to compute the stream of interest tax shields and its present value, $PV(ITS)$. What is the
relation between $PV(ITS)$, the value of the assets of the levered firm and results of part (b)?\\
\emph{Hint: Interest tax shields at t=1 are equal to $D_0\times r_D\times\tau$, where $D_0$ is market
value of debt at $t=0$; interest tax shields at t=2 are equal to $D_1\times r_D\times\tau$, where $D_1$
is market value of debt at $t=1$; and so on.}}
